\documentclass[11pt,letterpaper]{article}

\usepackage[T1]{fontenc}
\usepackage{lmodern}
\usepackage{microtype}
\usepackage[margin=1.05in]{geometry}
\usepackage{amsmath}
\usepackage{booktabs}
\usepackage{xcolor}
\usepackage{listings}
\usepackage{hyperref}

\definecolor{accent}{HTML}{174D61}
\definecolor{codebackground}{HTML}{F2F5F6}
\hypersetup{
  colorlinks=true,
  linkcolor=accent,
  urlcolor=accent,
  pdftitle={Fundamentals of Numerical Computation - Study Notes}
}
\lstset{
  language=Python,
  basicstyle=\small\ttfamily,
  backgroundcolor=\color{codebackground},
  keywordstyle=\color{accent}\bfseries,
  showstringspaces=false,
  columns=fullflexible,
  keepspaces=true,
  xleftmargin=0.5em,
  xrightmargin=0.5em,
  frame=single,
  rulecolor=\color{codebackground}
}

\newcommand{\reading}[2]{\par\smallskip
  {\small Textbook reading \enspace\href{#1}{#2}}\par\medskip}
\DeclareMathOperator{\fl}{fl}
\setlength{\parindent}{0pt}
\setlength{\parskip}{0.45em}
\setcounter{tocdepth}{2}

\begin{document}

{\color{accent}\large Numerical Analysis Study Notes\par}
\vspace{0.45em}
{\LARGE\bfseries Fundamentals of Numerical Computation\par}
\smallskip
{\large Preface and Chapter 1\par}
\vspace{1em}
An independent reading companion to Tobin A.\ Driscoll and Richard J.\
Braun's \href{https://fncbook.com/python/}{Python edition}.

\section*{Preface}
\addcontentsline{toc}{section}{Preface}
Numerical computation begins with a mathematical question and ends with an
answer produced by finite arithmetic. This pilot studies what happens
between those points. The examples use calculus to predict sensitivity and
Python to check how arithmetic and algorithm choices affect the result.

As you work through a section, write down the exact mathematical target,
predict where accuracy could be lost, and run an experiment against a
reference value. Record the input and the error measure alongside the
result. The examples here use only the Python standard library. For the
book's own examples, see its setup instructions.

{\small
  \href{https://fncbook.com/python/preface-original/}{Original preface}
  \enspace\textbullet\enspace
  \href{https://fncbook.com/python/preface/}{Julia edition preface}
  \enspace\textbullet\enspace
  \href{https://fncbook.com/python/usage/}{How to use the book}
  \enspace\textbullet\enspace
  \href{https://fncbook.com/python/setup/}{Python setup}\par
}

\tableofcontents

\clearpage
\section{Introduction}
\reading{https://fncbook.com/python/overview/}{Chapter overview}

Roundoff can change the inputs and the intermediate results of a computation.
Conditioning describes how sensitive the exact answer is to changes in the
inputs. Stability describes how the chosen algorithm handles the errors
introduced during computation.

\subsection{Floating-point numbers}
\reading{https://fncbook.com/python/floating-point/}{Section 1.1}

For IEEE binary64, a finite, nonzero \emph{normal} number has the form
\[
  \pm 2^e \left(1+\frac{k}{2^{52}}\right),
  \qquad
  -1022 \le e \le 1023,\qquad 0 \le k < 2^{52},
\]
where $e$ and $k$ are integers. Within $[2^e,2^{e+1})$ the spacing is
$2^{e-52}$. In particular, the distance from $1$ to the next larger
representable number is $\epsilon_{\rm mach}=2^{-52}$. Under rounding to
nearest, the \emph{unit roundoff} is $u=\epsilon_{\rm mach}/2=2^{-53}$.
For a real $x$ in the normal range with a finite rounded value,
\[
  \fl(x)=x(1+\delta),\qquad |\delta|\le u.
\]
Subnormal numbers have fixed absolute spacing and may exceed this relative
error bound. Overflow can produce an infinity.

For an approximation $\widehat{x}$ to a target $x$, absolute error is
$|\widehat{x}-x|$. Relative error divides that quantity by $|x|$, so the
ordinary relative error is defined only for $x\ne0$.

\paragraph{Experiment.} A spacing of $4$ at $2^{54}$ explains why adding
$1$ there rounds back to the original number. Reordering the same terms
can change the result.

\begin{lstlisting}
from math import nextafter, inf

x = float(2**54)
print(nextafter(x, inf) - x)  # 4.0
print((x + 1.0) - x)         # 0.0
print((x - x) + 1.0)         # 1.0
\end{lstlisting}

\textbf{Check your understanding.} Predict which of $2^{54}+2$ and
$2^{54}+3$ round up under ties-to-even. Verify by subtracting $2^{54}$ from
the floating-point result.

\clearpage
\subsection{Problems and conditioning}
\reading{https://fncbook.com/python/conditioning/}{Section 1.2}

Suppose the desired output is $y=f(x)$ and the input is perturbed by $h$.
Taylor expansion gives $f(x+h)-f(x)=f'(x)h+o(h)$ for a differentiable
scalar $f$. If $x$ and $f(x)$ are nonzero, the ratio of relative output
change to relative input change tends to
\[
   \kappa_f(x)=\left|\frac{x f'(x)}{f(x)}\right|.
\]
Thus $\kappa_f(x)$ estimates how much a small relative input error can
grow in the exact problem. At a zero of $f$, use absolute output changes
or choose an output scale appropriate to the problem.

\paragraph{Example.} For $f(x)=1/(1-x)$, the relative condition number is
$|x/(1-x)|$. At $x=0.999$ it equals $999$. An input change of $10^{-6}$
raises the exact output from $1000$ to $1/(0.000999)\approx1001.001$.
The input changes by about one part in a million, and the output changes
by about one part in a thousand.

\subsection{Algorithms}
\reading{https://fncbook.com/python/algorithms/}{Section 1.3}

The evaluation order determines the operations an algorithm performs. For
example, write the polynomial $3x^3-2x^2+5x-7$ in nested form,
\[
   ((3x-2)x+5)x-7.
\]
For a degree-$n$ polynomial, Horner's rule starts with the leading
coefficient, then multiplies by $x$ and adds the next coefficient at
each step. It uses $n$ multiplications and $n$ additions when counted
separately.

\begin{lstlisting}
def evaluate_polynomial(coefficients, x):
    leading, *remaining = coefficients
    value = leading
    for coefficient in remaining:
        value = value * x + coefficient
    return value

assert evaluate_polynomial((3, -2, 5, -7), 2.0) == 19.0
\end{lstlisting}

The coefficients here run from highest degree to constant term. To check
the loop by hand, label the accumulator after each addition and expand the
result.

\clearpage
\subsection{Stability}
\reading{https://fncbook.com/python/stability/}{Section 1.4}

For a computed result $\widehat{y}$, the absolute \emph{forward error} is
$|\widehat{y}-f(x)|$. The absolute \emph{backward error} is the smallest
$|\Delta x|$ for which $f(x+\Delta x)=\widehat{y}$, when such a smallest
perturbation exists. A backward-stable algorithm produces the exact
answer for data that differ from the given input by a small amount on
the scale of the data. If the problem is well conditioned, a small
backward error also implies a small forward error locally.

\paragraph{Experiment.} For $x>0$ let $f(x)=\sqrt{x^2+1}-x$. Since
$f'(x)=-f(x)/\sqrt{x^2+1}$, its relative condition number is at most $1$.
As $x$ grows, the two terms in the subtraction get closer, and rounding
in the square root can dominate their difference. The identity
\[
   f(x)=\frac{1}{\sqrt{x^2+1}+x}
\]
gives a second evaluation order.

\begin{lstlisting}
from math import hypot

for x in (1.0, 1e4, 1e8):
    root = hypot(x, 1.0)
    direct = root - x
    rationalized = 1.0 / (root + x)
    print(f"{x:8.1e}  {direct:.12e}  {rationalized:.12e}")
\end{lstlisting}

\begin{center}
  \small
  \begin{tabular}{@{}rcc@{}}
    \toprule
    $x$ & Subtraction & Rationalized form \\
    \midrule
    $1$      & $4.142135623731\times10^{-1}$ & $4.142135623731\times10^{-1}$ \\
    $10^4$   & $5.000000055588\times10^{-5}$ & $4.999999987500\times10^{-5}$ \\
    $10^8$   & $0$ & $5.000000000000\times10^{-9}$ \\
    \bottomrule
  \end{tabular}
\end{center}

At $x=10^8$, the subtraction gives zero. Since $f$ is strictly positive
for every finite real input, no finite input perturbation explains this
computed value. Use \texttt{decimal} at higher precision to estimate
the forward error of each method at the three test inputs.

\subsection*{Next steps}
\addcontentsline{toc}{subsection}{Next steps}
In Chapter 2, apply the same distinction to linear systems. For a computed
solution $\widehat{x}$ to $Ax=b$, first inspect the residual
$r=b-A\widehat{x}$. For a nonsingular $A$, the exact solution error
satisfies $x-\widehat{x}=A^{-1}r$. The action of $A^{-1}$ can turn a
small residual into a much larger error.

{\small
  \href{https://fncbook.com/next/}{Chapter 1 further reading}
  \enspace\textbullet\enspace
  \href{https://fncbook.com/python/overview-1/}{Chapter 2 overview}\par
}

\end{document}
